% Recovery of 09c_entropia_sectorial.tex, 11_confluencia.tex, and
% 85_agujero_negro_doble_hoja.tex; originals in fuentes_conservadas/
% horizontes_informacion. Provenance: desarrollo/HORIZONTES_INFORMACION_FUENTES.md.
% References to antecedents are agreed with the editor of VII.
% Does not incorporate a material current or evaluate the torsional amplitude.

\section{Pure bipartite states, reduced entropy, and horizon area laws}
\label{vii:sec:horizontes-informacion}

The enriched state retains paths and their registers during
evolution; a subsystem readout determines which part of that information
remains accessible. The construction in this section uses the angular
channels of Section~\ref{iv:angular:section}, the incidence and
area operator of Section~\ref{v:ant:sec:area}, and the units
of action, area, and temperature of
Sections~\ref{v:ant:sec:gravity}
and~\ref{v:ant:sec:ii-kb-aut-desarrollo}. These are earlier representations of the
same APP--TRIT--TPK construction. Each transmits its domain,
orientation, normalization, and realization conditions here.

Three spaces are distinguished: the binary sheet register, the
tritic factors of sectorial incidence, and the geometric realization of the
exterior regions of a horizon. Purification, transport of observables,
and dimensional conversion specify the relations between them. This
distinction preserves the structure accompanying each entropy value.

\subsection{Sheet register and bipartite purification}
\label{vii:subsec:hojas-purificacion}

Write \(a=A_{\rm rad}\) and \(y=C^*_{\rm rad}\) for the
angular coordinates already constructed. The channels and their probabilities are
\begin{equation}
 \mathfrak q_\pm=e^{-a\mp y},\qquad
 p_\pm=\frac{\mathfrak q_\pm}{\mathfrak q_++\mathfrak q_-}
       =\frac{e^{\mp y}}{2\cosh y}.
 \label{vii:eq:probabilidades-hoja}
\end{equation}
The signs index these two channels. On
\(\mathcal H_{\rm h}=\operatorname{span}\{|+\rangle,|-\rangle\}\)
define
\[
 \rho_{\rm h}(y)=p_+|+\rangle\langle+|+p_-|-\rangle\langle-|.
\]
The auxiliary space
\(\mathcal H_{\rm ar}=\operatorname{span}
\{|\mathrm{adv}\rangle,|\mathrm{ret}\rangle\}\)
carries two orthonormal labels. With a phase \(\chi_0\) fixed, the joint
state is
\begin{equation}
 |\Psi_{\rm ar}\rangle
 =\sqrt{p_+}|+\rangle\otimes|\mathrm{adv}\rangle
  +e^{i\chi_0}\sqrt{p_-}|-\rangle\otimes|\mathrm{ret}\rangle.
 \label{vii:eq:purificacion-hoja}
\end{equation}

\begin{proposition}[Accessible information in the sheet register]
\label{vii:prop:purificacion-hoja}
The vector~\eqref{vii:eq:purificacion-hoja} is normalized, and its partial
trace over \(\mathcal H_{\rm ar}\) is \(\rho_{\rm h}(y)\).
Its entanglement entropy, expressed in nats, is
\begin{equation}
 s_{\rm h}(y)=\log(2\cosh y)-y\tanh y.
 \label{vii:eq:entropia-hoja}
\end{equation}
Exchanging the labels \(+\) and \(-\) conjugates
\(\rho_{\rm h}(y)\) to \(\rho_{\rm h}(-y)\) and preserves
\(s_{\rm h}\).
\end{proposition}
\begin{proof}
The squared norm is \(p_++p_-=1\). Orthogonality of the
auxiliary labels eliminates the cross terms in the partial trace.
The squared Schmidt coefficients are precisely \(p_\pm\).
By~\eqref{vii:eq:probabilidades-hoja},
\(p_+-p_-=-\tanh y\); substituting
\(\log p_\pm=\mp y-\log(2\cosh y)\) into
\(-\sum_\pm p_\pm\log p_\pm\) proves the formula.
Exchanging the signs permutes the two eigenvalues and preserves entropy.
\end{proof}

The common factor \(e^{-a}\) cancels when the binary register is normalized.
The relative phase \(\chi_0\) remains in the joint state even though this
diagonal reduction does not express it. Logical reversibility concerns
transport of the complete register, with its labels and correlations;
phase return preserves the memory advance described in
Section~\ref{vii:sec:retorno-memoria-gravedad}. A thermal implementation
also has its own energy and dissipation law. In particular,
\(s_{\rm h}\) measures this binary subsystem and its specific bipartition.

\subsection{Purification of the incidence sectors}
\label{vii:subsec:purificacion-sectorial}

The boundary incidence constructed earlier distinguishes two sets
of eighteen links, with sectorial labels \(m=90,120\). Each link
carries \(\mathbb C^3\), with basis \(|0\rangle,|1\rangle,|2\rangle\)
and occupation operator \(\operatorname{diag}(0,1,2)\). Let
\(\mathsf N_m\) be their sectorial sums, and
\begin{equation}
 \mathsf N=\mathsf N_{90}+\mathsf N_{120},\qquad
 \mathsf D_\partial=90\mathsf N_{90}+120\mathsf N_{120}.
 \label{vii:eq:ocupaciones-sectoriales}
\end{equation}
On the positive branch, \(\gamma_{\rm BI}\) denotes the Barbero--Immirzi
parameter already obtained, and \(a_0>0\) its angular area factor.
This notation distinguishes the parameter from the connection and its holonomies.
Area is expressed as
\begin{equation}
 \widehat{\mathcal A}_m=\ell_P^2\gamma_{\rm BI}a_0\mathsf N_m,
 \qquad
 \widehat{\mathcal A}=\widehat{\mathcal A}_{90}
                         +\widehat{\mathcal A}_{120}.
 \label{vii:eq:area-sectores}
\end{equation}

For the finite parameter \(\Delta=\Delta_{\rm mod}\in\mathbb R\) of the
reduced state, put
\begin{align}
 x_m&=e^{-m\Delta},& z_m&=1+x_m+x_m^2,&
 Z(\Delta)&=z_{90}^{18}z_{120}^{18},\\
 \rho_\partial(\Delta)&=Z(\Delta)^{-1}e^{-\Delta\mathsf D_\partial}.
 \label{vii:eq:estado-sectorial}
\end{align}
The trace in this section is the ordinary matrix trace, and the states
have trace one. The parameter \(\Delta\) retains the selection and
normalization of the antecedent state.

\begin{proposition}[Tensor-product purification of the boundary]
\label{vii:prop:purificacion-sectorial}
For a finite real weight \(w\), let
\begin{equation}
 |\operatorname{TFD}(w)\rangle
 =\frac{|0\rangle|0\rangle+e^{-w/2}|1\rangle|1\rangle
                   +e^{-w}|2\rangle|2\rangle}
        {\sqrt{1+e^{-w}+e^{-2w}}}.
 \label{vii:eq:tfd-local}
\end{equation}
The product
\begin{equation}
 |\Psi_\Delta\rangle
 =|\operatorname{TFD}(90\Delta)\rangle^{\otimes18}
  \otimes|\operatorname{TFD}(120\Delta)\rangle^{\otimes18}
 \label{vii:eq:tfd-sectorial}
\end{equation}
is a normalized purification of \(\rho_\partial(\Delta)\).
\end{proposition}
\begin{proof}
At each link, the sum of the squares of the three coefficients is one.
The partial trace eliminates terms with distinct indices and returns
\(\operatorname{diag}(1,e^{-w},e^{-2w})/(1+e^{-w}+e^{-2w})\).
The partial trace commutes with the tensor product. Multiplying the
thirty-six reductions gives the exponential and the partition function
in~\eqref{vii:eq:estado-sectorial}.
\end{proof}

\subsection{Information equation and area fluctuations}
\label{vii:subsec:ecuacion-informacional}

The occupation means and variances per link are
\begin{equation}
 f_m=\frac{x_m+2x_m^2}{z_m},\qquad
 v_m=\frac{x_m+4x_m^2}{z_m}-f_m^2.
 \label{vii:eq:momentos-sectoriales}
\end{equation}
Write \(s(\rho)=-\operatorname{Tr}(\rho\log\rho)\), with
\(0\log0=0\). In the preceding purification, \(s(\rho_\partial)\)
is its entanglement entropy.

\begin{proposition}[Mean area, entropy, and sectorial response]
\label{vii:prop:respuesta-sectorial}
For real, finite \(\Delta\), we have
\begin{align}
 \langle\mathsf N\rangle_\Delta&=18(f_{90}+f_{120}),\\
 \langle\mathsf D_\partial\rangle_\Delta&=18(90f_{90}+120f_{120}),\\
 \langle\widehat{\mathcal A}\rangle_\Delta
 &=18\ell_P^2\gamma_{\rm BI}a_0(f_{90}+f_{120}),\\
 s(\Delta)&=\log Z(\Delta)+\Delta\langle\mathsf D_\partial\rangle_\Delta,
 \label{vii:eq:entropia-sectorial}\\
 \frac{\mathrm d\langle\mathsf D_\partial\rangle_\Delta}{\mathrm d\Delta}
 &=-\operatorname{Var}_\Delta(\mathsf D_\partial)
  =-18(90^2v_{90}+120^2v_{120}),\\
 \frac{\mathrm ds}{\mathrm d\Delta}
 &=-\Delta\operatorname{Var}_\Delta(\mathsf D_\partial),\\
 \frac{\mathrm d\langle\widehat{\mathcal A}\rangle_\Delta}{\mathrm d\Delta}
 &=-18\ell_P^2\gamma_{\rm BI}a_0(90v_{90}+120v_{120}).
 \label{vii:eq:respuestas-area-entropia}
\end{align}
\end{proposition}
\begin{proof}
The local weights are \((1,x_m,x_m^2)/z_m\), whose first two
moments give \(f_m\) and \(v_m\). Additivity of the occupation
operators and factorization of the state give the first three formulas.
Taking the trace of
\(-\rho_\partial\log\rho_\partial
=\rho_\partial(\log Z+\Delta\mathsf D_\partial)\)
gives entropy. Differentiation of the finite sums yields
\(\partial_\Delta\log Z=-\langle\mathsf D_\partial\rangle_\Delta\)
and \(\partial_\Delta^2\log Z
=\operatorname{Var}_\Delta(\mathsf D_\partial)\).
The independent factors add their variances. In the derivative of
\eqref{vii:eq:entropia-sectorial}, the two first-moment terms
cancel. Finally, \(\partial_\Delta f_m=-mv_m\)
gives the area response.
\end{proof}

Loss of uniformity is quantified in dimension \(d_\partial=3^{36}\):
\begin{equation}
 I_\partial(\Delta)
 =36\log3-s(\Delta)
 =\mathcal D\bigl(\rho_\partial(\Delta)\Vert I/d_\partial\bigr)\ge0.
 \label{vii:eq:deficit-informacional}
\end{equation}
Here \(\mathcal D\) denotes relative entropy, defined and justified
below. Substituting \(\log(I/d_\partial)=-36\log3\,I\) proves the
identity. Convexity of \(t\log t\) gives positivity; equality
corresponds to the uniform distribution, which in this collection occurs
exactly at \(\Delta=0\). The deficit measures dimensionless information;
the scale \(\gamma_{\rm BI}a_0\ell_P^2\) measures area per occupation.

\begin{theorem}[Occupation reversal and complementary area]
\label{vii:thm:inversion-sectorial}
Let \(\mathcal R_\partial\) be the involution exchanging
\(|0\rangle\) and \(|2\rangle\) and fixing \(|1\rangle\) in each
factor. Then
\begin{align}
 \mathcal R_\partial\mathsf N\mathcal R_\partial^*&=72I-\mathsf N,\\
 \mathcal R_\partial\mathsf D_\partial\mathcal R_\partial^*
 &=7560I-\mathsf D_\partial,\\
 \rho_\partial(-\Delta)&=\mathcal R_\partial\rho_\partial(\Delta)
                          \mathcal R_\partial^*,\\
 s(-\Delta)&=s(\Delta),\\
 \langle\widehat{\mathcal A}\rangle_{-\Delta}
 +\langle\widehat{\mathcal A}\rangle_\Delta
 &=72\gamma_{\rm BI}a_0\ell_P^2,\\
 \mathcal D\bigl(\rho_\partial(\Delta)\Vert\rho_\partial(-\Delta)\bigr)
 &=\Delta\bigl(7560-2\langle\mathsf D_\partial\rangle_\Delta\bigr).
 \label{vii:eq:inversion-sectorial}
\end{align}
\end{theorem}
\begin{proof}
Each occupation transforms by \(n\mapsto2-n\). In each sector,
\(\mathsf N_m\mapsto36I-\mathsf N_m\), so the constant
term of the weighted reader is \(36(90+120)=7560\).
Moreover, \(z_m(-\Delta)=e^{2m\Delta}z_m(\Delta)\) and
\(Z(-\Delta)=e^{7560\Delta}Z(\Delta)\). Substitution into the
normalized state proves conjugacy. Its spectrum preserves entropy;
the transformation of \(\mathsf N\) gives the complementary area. Subtracting
the logarithms of the two states and taking the expectation in
\(\rho_\partial(\Delta)\) gives the last equality.
\end{proof}

At \(\Delta=0\), \(s=36\log3\). As \(\Delta\to+\infty\),
the state concentrates on zero occupation and entropy tends to zero;
reversal takes this limit to maximum occupation. For \(\Delta>0\),
entropy decreases strictly because the weighted reader has more
than one eigenvalue and the state has full rank. The complete
configuration, area, and entropy thus retain distinct roles.

\subsection{Modular law and finite area variations}
\label{vii:subsec:ley-modular-finita}

Fix \(\rho_0=\rho_\partial(\Delta_0)\), with
real, finite \(\Delta_0\), and hold
\(\Delta_0,\gamma_{\rm BI},a_0,\ell_P\) constant. The normalized
areas \(\mathcal A_m^{\rm n}=\gamma_{\rm BI}a_0\mathsf N_m\)
express the modular operator as
\begin{equation}
 \mathsf K_0=-\log\rho_0
 =c_0I+\sum_{m=90,120}\beta_m\mathcal A_m^{\rm n},
 \quad \beta_m=\frac{m\Delta_0}{\gamma_{\rm BI}a_0},
 \quad c_0=\log Z(\Delta_0).
 \label{vii:eq:operador-modular-areal}
\end{equation}
We have \(\beta_{120}=\tfrac43\beta_{90}\), including when both
are zero. The ratio is \(4/3\) when \(\Delta_0\ne0\).
\(\mathsf K_0\) is the logarithm of this reference state; its type
distinguishes it from temporal transport and the gravitational connection.

\begin{theorem}[First modular law and finite correction]
\label{vii:thm:ley-modular-finita}
For any normalized state \(\sigma\) on the same space, let
\(\Delta_\sigma\langle B\rangle
=\operatorname{Tr}[(\sigma-\rho_0)B]\). Then
\begin{equation}
 s(\sigma)-s(\rho_0)
 =\sum_{m=90,120}\beta_m\Delta_\sigma\langle\mathcal A_m^{\rm n}\rangle
  -\mathcal D(\sigma\Vert\rho_0),
 \label{vii:eq:ley-modular-finita}
\end{equation}
where
\(\mathcal D(\sigma\Vert\rho_0)
=\operatorname{Tr}[\sigma(\log\sigma-\log\rho_0)]\)
is finite and nonnegative. If \(\sigma(\varepsilon)\) is differentiable,
normalized, and \(\sigma(0)=\rho_0\), we obtain
\begin{equation}
 \left.\frac{\mathrm ds(\sigma(\varepsilon))}{\mathrm d\varepsilon}\right|_0
 =\sum_{m=90,120}\beta_m
   \left.\frac{\mathrm d\langle\mathcal A_m^{\rm n}\rangle_{
                 \sigma(\varepsilon)}}{\mathrm d\varepsilon}\right|_0.
 \label{vii:eq:primera-ley-modular}
\end{equation}
\end{theorem}
\begin{proof}
By definition,
\(\mathcal D(\sigma\Vert\rho_0)=-s(\sigma)
+\operatorname{Tr}(\sigma\mathsf K_0)\), whereas
\(s(\rho_0)=\operatorname{Tr}(\rho_0\mathsf K_0)\).
Subtraction and~\eqref{vii:eq:operador-modular-areal} prove the finite
identity, since normalization cancels the term \(c_0I\).

For positivity, diagonalize \(\sigma\) with eigenvalues
\(p_i\) and vectors \(u_i\), and \(\rho_0\) with eigenvalues
\(q_j>0\) and vectors \(v_j\). If
\(r_i=\langle u_i,\rho_0u_i\rangle>0\), concavity of the
logarithm gives
\[
 \sum_j|\langle u_i,v_j\rangle|^2\log q_j\le\log r_i.
\]
Consequently,
\(\mathcal D(\sigma\Vert\rho_0)
\ge\sum_i p_i\log(p_i/r_i)\ge0\).
The last inequality is convexity of \(t\log t\) with weights
\(r_i\), whose sum is one. Faithfulness of \(\rho_0\) keeps
its logarithms finite; zero eigenvalues of \(\sigma\) are
interpreted through \(0\log0=0\).

For the infinitesimal formula, near the faithful state \(\rho_0\),
the derivative of the trace of \(-\sigma\log\sigma\) is
\(-\operatorname{Tr}[\dot\sigma(\log\rho_0+I)]\).
This identity follows from the spectral derivative of the trace.
The condition \(\operatorname{Tr}\dot\sigma=0\) leaves
\(\operatorname{Tr}(\dot\sigma\mathsf K_0)\), and substituting
\eqref{vii:eq:operador-modular-areal} completes the proof.
\end{proof}

For physical areas, the coefficients are \(\beta_m/\ell_P^2\).
Applying the common transducer \(k_B\) gives \(S=k_Bs\) and
retains the full term \(k_B\mathcal D\). The derivative in
\eqref{vii:eq:respuestas-area-entropia} varies \(\Delta\); law
\eqref{vii:eq:primera-ley-modular} holds \(\Delta_0\) fixed and
varies the state around the reference. These are different operations.

\subsection{Action, decision energy, and the area unit}
\label{vii:subsec:conversion-informacion-area}

Fix a common oriented action section and its gravitational and
thermal sections, on the positive branch. The antecedents provide
\begin{equation}
 \begin{aligned}
 \ell_P^2&=\frac{\hbar G}{c^3},& t_P&=\frac{\ell_P}{c},\\
 T_{108}&=2\pi t_P,& h&=2\pi\hbar,\\
 E_P&=\frac{\hbar}{t_P}=k_B\Theta_{\rm clk}\log3.&&
 \end{aligned}
 \label{vii:eq:unidad-accion-informacion}
\end{equation}
The thermal equality retains the pair and normalization previously
constructed; when their dimensional basis changes,
\(k_B\) and \(\Theta_{\rm clk}\) are transported together. The elementary area increment is
\(\delta\mathcal A=\gamma_{\rm BI}a_0\ell_P^2\).

\begin{proposition}[Conversion coefficient between decision and area]
\label{vii:prop:tension-informacion-area}
The ratio of decision energy to the area increment satisfies
\begin{equation}
 \begin{aligned}
 \mathcal T_{I/A}:=\frac{E_P}{\delta\mathcal A}
 &=\frac{k_B\Theta_{\rm clk}\log3}{\gamma_{\rm BI}a_0\ell_P^2}\\
 &=\frac{c^3}{\gamma_{\rm BI}a_0G t_P}
  =\frac{c^4}{\gamma_{\rm BI}a_0G\ell_P}.
 \end{aligned}
 \label{vii:eq:tension-informacion-area}
\end{equation}
Its dimension is energy per area.
\end{proposition}
\begin{proof}
Substitution of \(E_P=\hbar/t_P\) and
\(\ell_P^2=\hbar G/c^3\) cancels the common action section.
The identity \(\ell_P=ct_P\) gives the final expression. Positivity
of the factors ensures that every quotient is defined.
\end{proof}

The capacity \(81\) of the cubic cut and the bipartition of the thirty-six
sectorial links correspond to different domains. In the
cubic graph \(\{0,\ldots,N-1\}^3\), the \(N^2\) parallel lines
between two opposite faces are edge-disjoint: every cut
intersects all \(N^2\), and a transverse layer attains this bound.
For \(N=9\), the uniform product state of the trits in that layer
has \(s_{\rm corte}=81\log3\). Assigning one decision of the
same cycle to each link expresses
\begin{equation}
 S_{\rm corte}^{\rm fis}=81k_B\log3,\qquad
 E_{\rm corte}^{(1)}=81E_P
                  =\Theta_{\rm clk}S_{\rm corte}^{\rm fis}.
 \label{vii:eq:corte-energia-informacion}
\end{equation}
The proof combines the minimum cut, additivity of the product-state
entropy, and~\eqref{vii:eq:unidad-accion-informacion}. For the sectorial
state, the relevant finite equation remains
\eqref{vii:eq:ley-modular-finita}, with its relative entropy.

\subsection{Cosmological normalization of the horizon}
\label{vii:subsec:balance-horizonte}

For an areal radius \(R>0\), the cosmological realization considered
uses the triad
\begin{equation}
 E_\Lambda(R)=\frac{c^4R}{2G},\qquad
 S_{\rm H}(R)=\frac{k_B\pi R^2}{\ell_P^2},\qquad
 T_{\rm cos}(R)=\frac{\hbar c}{2\pi k_BR}.
 \label{vii:eq:horizonte-cosmologico}
\end{equation}
The factor \(2\pi\) belongs to this temperature normalization.
The area \(\mathcal A_{\rm H}=4\pi R^2\) expresses the same entropy
as \(S_{\rm H}/k_B=\mathcal A_{\rm H}/(4\ell_P^2)\).

\begin{theorem}[Gravitational cell of half action]
\label{vii:thm:celda-media-accion}
In~\eqref{vii:eq:horizonte-cosmologico},
\begin{align}
 T_{\rm cos}(R)S_{\rm H}(R)&=E_\Lambda(R),\\
 T_{\rm cos}(\ell_P)S_{\rm H}(\ell_P)
 &=\frac{E_P}{2}=\frac{k_B\Theta_{\rm clk}\log3}{2},\\
 E_\Lambda(\ell_P)t_P&=\frac{\hbar}{2},&
 E_\Lambda(\ell_P)T_{108}&=\frac h2.
 \label{vii:eq:media-accion-horizonte}
\end{align}
Holding the chart constants fixed and varying \(R\),
\begin{equation}
 T_{\rm cos}\,\mathrm dS_{\rm H}=2\,\mathrm dE_\Lambda,
 \qquad S_{\rm H}\,\mathrm dT_{\rm cos}=-\mathrm dE_\Lambda.
 \label{vii:eq:diferencial-horizonte}
\end{equation}
\end{theorem}
\begin{proof}
Multiplying temperature and entropy gives
\(\hbar cR/(2\ell_P^2)=c^4R/(2G)\).
At \(R=\ell_P\), the result is
\(\hbar/(2t_P)=E_P/2\). The two durations in
\eqref{vii:eq:unidad-accion-informacion} give the action relations.
For the differentials,
\(\mathrm dE_\Lambda=c^4\mathrm dR/(2G)\),
\(\mathrm dS_{\rm H}=2\pi k_BR\mathrm dR/\ell_P^2\), and
\(\mathrm dT_{\rm cos}=-T_{\rm cos}\mathrm dR/R\).
Substitution proves both identities.
\end{proof}

The product \(E_\Lambda=T_{\rm cos}S_{\rm H}\) expresses a
state identity. Its radial variations retain both terms
of~\eqref{vii:eq:diferencial-horizonte}; a complete thermodynamic interpretation
also specifies the corresponding temporal generator, work, and
material realization. In this chart, action and
gravitation fix the unit \(\ell_P^2\), while Boltzmann converts
information into dimensional entropy. Transport of sections with
fixed \(G\hbar\) preserves that unit when the section \(c\) is
held fixed, and preserves \(S_{\rm H}/k_B\) at fixed radius.

In the cell \(R=\ell_P\), \(S_{\rm H}/k_B=\pi\), so
the effective multiplicity is \(\Omega_{\rm eff}=e^\pi\).
For the tritic hyperbolic involution \(H_{\rm trit}\), with spectrum
\(\{1,-1\}\), diagonalization gives
\(\operatorname{spec}(e^{\pi H_{\rm trit}})=\{e^\pi,e^{-\pi}\}\).
This spectral relation specifies the expanding value; a microscopic
identification of both spaces additionally retains their realization
map and state.

\subsection{Kruskal chart and exchange of exterior regions}
\label{vii:subsec:kruskal-dos-hojas}

Fix a expressed state \(x\), its radius \(R=R(x)>0\), and the
positive sections \(c,G\). In the following chart, \(x,R,c,G\)
are fixed parameters; \(t,r\) are the spacetime coordinates.
The exterior realization of the two sheets is specified by
\begin{equation}
 f(r)=1-\frac Rr,\qquad
 \mathrm ds^2=-f(r)c^2\mathrm dt^2+f(r)^{-1}\mathrm dr^2
                 +r^2\mathrm d\Omega_2^2,\qquad r>R.
 \label{vii:eq:metrica-exterior}
\end{equation}
For the sheet reader \(\sigma\in\{+1,-1\}\), define
\begin{align}
 r_*&=r+R\log\left(\frac rR-1\right),\\
 \mathcal U&=-\sigma\exp\left(-\frac{ct-r_*}{2R}\right),&
 \mathcal V&=\sigma\exp\left(\frac{ct+r_*}{2R}\right).
 \label{vii:eq:carta-kruskal}
\end{align}
The symbols \(\mathcal U,\mathcal V\) are coordinates and are
distinguished from path-transport operators.

\begin{theorem}[Geometric realization and sheet involution]
\label{vii:thm:kruskal-hojas}
The chart~\eqref{vii:eq:carta-kruskal} realizes the two sheets as the
two exterior regions \(\mathcal U\mathcal V<0\) of the Kruskal
extension. We have
\begin{equation}
 -\mathcal U\mathcal V=\left(\frac rR-1\right)e^{r/R},\qquad
 \mathrm ds^2=-\frac{4R^3}{r}e^{-r/R}\,
                    \mathrm d\mathcal U\,\mathrm d\mathcal V
                    +r^2\mathrm d\Omega_2^2.
 \label{vii:eq:metrica-kruskal}
\end{equation}
The involution
\begin{equation}
 \mathcal C_{\rm H}:(\sigma,\mathcal U,\mathcal V)
 \longmapsto(-\sigma,-\mathcal U,-\mathcal V)
 \label{vii:eq:involucion-horizonte}
\end{equation}
is isometric, fixes the areal radius, exchanges the exterior regions, and preserves
the set of horizons \(\mathcal U\mathcal V=0\).
The metric is regular at \(r=R\) and satisfies the vacuum equations
on the regular domain \(r>0\).
\end{theorem}
\begin{proof}
Multiplying the coordinates gives the first identity. Since
\(\mathrm dr_*/\mathrm dr=f^{-1}\),
\[
 \mathrm d\mathcal U=\frac{\mathcal U}{2R}
            (-c\,\mathrm dt+f^{-1}\mathrm dr),\qquad
 \mathrm d\mathcal V=\frac{\mathcal V}{2R}
            ( c\,\mathrm dt+f^{-1}\mathrm dr).
\]
Their product reproduces the radial and temporal part of the metric.
The function \(z\mapsto(z-1)e^z\) has derivative \(ze^z>0\)
for \(z>0\), which determines \(r\) from the product.
In the exterior one recovers
\(\sigma=\operatorname{sign}\mathcal V\) and
\(t=(R/c)\log(-\mathcal V/\mathcal U)\).
At \(r=R\), the radial coefficient of the extended metric is
\(-4R^2/e\), finite and nonzero. Reversal preserves the
product, radius, and \(\mathrm d\mathcal U\mathrm d\mathcal V\),
and exchanges the two sign choices.

To check the vacuum equations in the exterior chart, the independent mixed
components of the Einstein tensor reduce to
\[
 \mathsf G^t{}_t=\mathsf G^r{}_r
 =\frac{rf'+f-1}{r^2},\qquad
 \mathsf G^\theta{}_\theta=\mathsf G^\phi{}_\phi
 =\frac{f'}r+\frac{f''}2.
\]
With \(f'=R/r^2\) and \(f''=-2R/r^3\), all vanish.
The regular expression~\eqref{vii:eq:metrica-kruskal} extends the
result across the horizons on the domain \(r>0\).
\end{proof}

This realization preserves
\(E_{\rm H}=c^4R/(2G)\) and
\(S_{\rm H}=k_B\pi R^2/\ell_P^2\), which depend on the areal radius.
The temperature \(T_{\rm cos}\) in
\eqref{vii:eq:horizonte-cosmologico} belongs to the cosmological
realization declared there. The exterior chart
\eqref{vii:eq:metrica-exterior} retains its own temporal normalization.
The geometric map establishes the exchange of exterior regions; comparison
with collapse formation, radiation spectra, or quasinormal modes
also uses their respective physical conditions.

\subsection{Transport of states, observables, and refinements}
\label{vii:subsec:transporte-informacion-horizonte}

Let \(F_x\) denote the chart of the two exterior regions just constructed.
In each regular angular chart, its temporal and radial Jacobian is
\begin{equation}
 \left|\frac{\partial(\mathcal U,\mathcal V)}{\partial(t,r)}\right|
 =\frac{cr}{2R^3}e^{r/R}>0.
 \label{vii:eq:jacobiano-kruskal}
\end{equation}
The identity follows from the two differentials in the preceding proof.
The state's remaining registers are transported as fibre
coordinates; the chart preserves their identity and memory.

\begin{proposition}[Unitary equivalence of sheet and exterior readouts]
\label{vii:prop:transporte-unitario-horizonte}
Let \(\mu\) be a sheet measure and \(\nu=(F_x)_*\mu\) its
transported measure. The change of variables defines a unitary
\begin{equation}
 W_{\rm H}:L^2(\mu)\longrightarrow L^2(\nu),\qquad
 (W_{\rm H}\psi)(z)=\psi(F_x^{-1}z).
 \label{vii:eq:unitaria-horizonte}
\end{equation}
With densities relative to fixed coordinate measures, the same
map incorporates the corresponding Jacobian factor.
The conjugation
\begin{equation}
 \mathfrak I_{\rm H}(B)=W_{\rm H}BW_{\rm H}^*
 \label{vii:eq:mapa-observables-horizonte}
\end{equation}
is a unital completely positive \(*\)-isomorphism. It transports the
state, its expectations, and its modular functional calculus. For the symmetric
measure on the two sheets, it also intertwines their involutions.
\end{proposition}
\begin{proof}
The definition of pushforward measure gives
\(\int|\psi\circ F_x^{-1}|^2\,\mathrm d\nu
=\int|\psi|^2\,\mathrm d\mu\).
The chart is invertible on the exterior regions, so the isometry is
surjective. Unitary conjugation preserves products, adjoints, and
the identity; at each matrix amplification it conjugates by
\(I_n\otimes W_{\rm H}\), proving complete positivity.
For a state \(\rho\), cyclicity of the trace gives
\(\operatorname{Tr}(W_{\rm H}\rho W_{\rm H}^*
\,\mathfrak I_{\rm H}(B))=\operatorname{Tr}(\rho B)\).
Functional calculus implies
\(g(W_{\rm H}\rho W_{\rm H}^*)=W_{\rm H}g(\rho)W_{\rm H}^*\)
on its domain, in particular for the logarithm and modular flow of
a state faithful on its support. The equality
\(F_x\circ\mathcal C_{\rm hoja}=\mathcal C_{\rm H}\circ F_x\)
proves the intertwining of involutions; the symmetric measure realizes them
unitarily.
\end{proof}

\pagebreak % REV03: keep the closing refinement and bipartition discussion together.
If refinements preserve \((\sigma,t,r,R)\) and use the same
isometry on the fibre registers, the inclusions
\(j_m^{\rm hoja}\) and \(j_m^{\rm ext}\) satisfy
\begin{equation}
 W_{{\rm H},m+1}j_m^{\rm hoja}
 =j_m^{\rm ext}W_{{\rm H},m}.
 \label{vii:eq:naturalidad-horizonte}
\end{equation}
Indeed, the change of variables acts on the preserved geometric coordinates
and the inclusion on the same fibre coordinates; evaluating
both compositions gives the same function and the same transported density.
The equality extends the map to the refined system without resetting the registers.

For a bipartition \(\mathcal H_A\otimes\mathcal H_B\), a
factorized transport \(U_A\otimes U_B\) satisfies
\[
 \rho'_A=U_A\rho_AU_A^*,\qquad s(\rho'_A)=s(\rho_A).
\]
The first identity follows from the partial trace and the second from the spectrum.
For a general unitary, the equivalent formulation also transports
the subsystem algebra:
\(\mathcal B(\mathcal H_A)\otimes I\mapsto
U(\mathcal B(\mathcal H_A)\otimes I)U^*\).
This preserves the same operational question about accessible
information. The entropy equalities apply to states with the relevant finite
entropies.

The relation between information and gravitation is expressed through
distinct, composable operations: the state and bipartition fix accessible
information; incidence and Barbero--Immirzi fix area
resolution; action and gravitation construct the unit \(\ell_P^2\);
Boltzmann realizes entropy dimensionally; the horizon chart
transports geometry and its observables. The finite laws retain
relative entropy, and the geometric realizations retain their measures,
temporal normalization, and domains.
